\documentclass[10pt,twocolumn]{article}

% Standard packages
\usepackage[utf8]{inputenc}
\usepackage[T1]{fontenc}
\usepackage{hyperref}
\usepackage{graphicx}
\usepackage{booktabs}
\usepackage{amsmath}
\usepackage{xcolor}
\usepackage[margin=1in]{geometry}
\usepackage{natbib}
\usepackage{url}
% Allow URLs in the bibliography to break at any alphanumeric character —
% without this, the long Milvus blog URL is one unbreakable box that
% overflows the page (and its hyperlink annotation extends past the
% MediaBox, which tectonic warns about).
\def\UrlBreaks{\do\/\do-\do.\do=\do_\do?\do\&\do a\do b\do c\do d\do e\do f\do g\do h\do i\do j\do k\do l\do m\do n\do o\do p\do q\do r\do s\do t\do u\do v\do w\do x\do y\do z\do 0\do 1\do 2\do 3\do 4\do 5\do 6\do 7\do 8\do 9}
\usepackage{listings}
\usepackage{float}

% Hyperref setup
\hypersetup{
  colorlinks=true,
  linkcolor=blue,
  citecolor=blue,
  urlcolor=blue
}

% Listings setup for code
\lstset{
  basicstyle=\ttfamily\small,
  breaklines=true,
  frame=single,
  columns=fullflexible
}

\title{\textbf{sqz: Pre-Injection Context Compression\\for Agentic Code Generation}}

\author{
  Ojus Chugh\\
  Independent Researcher\\
  \texttt{ojuschugh@gmail.com}\\
  ORCID: \href{https://orcid.org/0009-0003-2863-1808}{0009-0003-2863-1808}
}

\date{May 2026 (revised September 2026, reflecting sqz 1.8.0)}

\begin{document}

\maketitle

% ============================================================
\begin{abstract}
Large Language Model (LLM) coding agents process tens of thousands of tokens per session through tool outputs---shell commands, file reads, test results, and build logs. These outputs enter the context window raw, consuming budget that could serve reasoning. We present \textbf{sqz}, a transparent compression layer that operates \emph{before} tool outputs are injected into context. sqz combines domain-specific structural formatters, content-addressed deduplication, adaptive pressure-aware compression, and entropy-based safety routing to achieve a mean 24.7\% reduction on first-pass content and up to 92\% savings on repeated reads---without requiring model fine-tuning, prompt modification, or changes to the agent's workflow. Measured across 3,003 real compressions in production agentic sessions, sqz saved 178,442 tokens while maintaining full semantic fidelity on downstream tasks.
\end{abstract}

% ============================================================
\section{Introduction}

The emergence of agentic coding tools---Claude Code, Cursor, Codex, Kiro, Gemini CLI---has shifted LLM usage from single-turn question-answering to multi-turn, tool-augmented sessions that can span thousands of interactions. In these sessions, the dominant source of context consumption is not user prompts or model responses, but \textbf{tool outputs}: the results of shell commands, file reads, and API calls that accumulate in the context window over time.

Recent measurements show that average agent session length tripled from under 2,000 tokens in late 2023 to over 5,400 tokens by late 2025, with the bulk of that growth coming from tool results accumulating in context~\cite{pan2026tool}. This creates a triple penalty: higher costs, increased latency from longer prefills, and degraded reasoning quality as attention must span an ever-growing context.

Prior work on context compression has focused on two regimes: (1)~\textbf{prompt compression} that removes redundant tokens from instructions before inference~\cite{jiang2023llmlingua,pan2024llmlingua2,wang2024lightning}, and (2)~\textbf{post-hoc context pruning} that evicts stale conversation history after it enters the window~\cite{anagnostidis2023dynamic,acon2024}. Both operate on text \emph{already inside} the context. sqz addresses a third, underexplored regime: \textbf{pre-injection compression of tool outputs}---reducing token count before content ever enters the LLM's context window.

This distinction is architecturally significant. Pre-injection compression:
\begin{itemize}
  \item Preserves prompt cache validity (the message prefix never changes)
  \item Operates independently of the LLM provider
  \item Requires no model cooperation or fine-tuning
  \item Works transparently across all agent frameworks
\end{itemize}

% ============================================================
\section{System Architecture}

\begin{figure*}[t]
  \centering
  \includegraphics[width=0.85\textwidth]{sqz-architecture.png}
  \caption{sqz system architecture---integration surfaces, 8-stage compression pipeline, and supporting subsystems. The system operates with zero telemetry, fully offline, and is air-gap capable.}
  \label{fig:architecture}
\end{figure*}

sqz operates as a PreToolUse hook---a command that the agent framework invokes before delivering tool output to the model. The system comprises four integration surfaces (CLI Hook, MCP Server, Browser Extension, IDE Extension) feeding into a unified Rust core (\texttt{sqz\_engine}) with an 8-stage compression pipeline. Three supporting subsystems handle persistence, intelligence, and analytics (Figure~\ref{fig:architecture}).

The architecture has five logical stages:

\subsection{Domain-Specific Structural Formatters}

Rather than treating all text uniformly (as token-level compression methods do), sqz recognizes 90+ CLI output patterns and applies format-aware transformations:

\begin{itemize}
  \item \textbf{Test output} $\rightarrow$ failures only (passing tests stripped)
  \item \textbf{Git status} $\rightarrow$ compact change summary
  \item \textbf{JSON responses} $\rightarrow$ null-stripped, key-projected, array-collapsed
  \item \textbf{Docker/kubectl} $\rightarrow$ name/image/status table extraction
  \item \textbf{Build logs} $\rightarrow$ error-focused summarization
\end{itemize}

This exploits the observation from CodeComp~\cite{chen2025codecomp} that source code and structured outputs have formal semantics governed by structure rather than surface-level co-occurrence. A \texttt{cargo test} output of 500 lines can be faithfully represented by 30 lines containing only the failures and their locations.

\subsection{Content-Addressed Deduplication}

Agentic sessions exhibit high repetition: the same file is read multiple times, the same command re-run after edits, the same API response fetched across turns. sqz maintains a persistent SHA-256 content cache (surviving across process invocations via SQLite) and replaces repeated content with a 13-token reference marker (\texttt{\S ref:HASH\S}).

This is analogous to the demand paging concept described in~\cite{demandpaging2025}, where context content is loaded on-demand rather than kept resident. The key insight is that in coding sessions, \textbf{the dominant savings come not from compressing individual outputs harder, but from recognizing that the same content appears repeatedly.}

A reference is only emitted while the content it points at is plausibly still in the model's context: refs expire after a freshness window (30 minutes of wall-clock time by default) and are invalidated eagerly when the agent framework compacts its history, at which point the content is re-sent in full. As of 1.4.0, entries can also be \emph{pinned} (\texttt{sqz pin}), exempting them from the freshness check---useful for stable knowledge-base content (project documentation, system prompts) that multiple agents sharing the session store reference across days.

\textbf{Near-duplicate deltas.} A file re-read after a small edit is not byte-identical, so exact hashing misses it. sqz detects near-duplicates of cached entries and emits a line-level delta---only the changed lines, plus a reference to the version the model already has. Since 1.5.0 a session-lifetime MinHash LSH index selects delta candidates, so a file re-read late in a long session still matches its earlier version rather than only the most recent cache entries.

\textbf{Ranged (slice) references, 1.8.0.} A field measurement over one user's real Claude Code sessions (measurement script published~\cite{harrison2026repeats}; counters reset at each compaction boundary, edit/write confirmations excluded) found that byte-identical repeats were rare (0.3\% of result characters; 1 in 542 file reads) but 8.4\% of file reads were a \emph{line range} of a file already read in full---\texttt{sed -n '41,80p'} after \texttt{cat}, or a ranged read tool call. Exact hashing misses these entirely, and near-duplicate matching does too: a 40-line slice of a 500-line file has a Jaccard similarity near 0.1. sqz therefore checks \emph{containment}: when new output is, on line boundaries, a slice of a fresh full read whose served bytes also contain the slice, it emits \texttt{\S ref:HASH:L41-80\S} instead of the lines. The second condition matters---the slice must appear in what was actually delivered to the model, so a ranged reference can never point at lines a lossy stage dropped. Guards: at least 160 bytes and two lines, and the same freshness and compaction rules as whole-content refs. The reference is not a summary; it is a pointer to bytes the model already holds, and \texttt{sqz expand 'HASH:L41-80'} recovers exactly those lines on demand.

Measured empirically: dedup accounts for 92\% reduction on repeat reads, compared to 21--58\% from first-pass structural compression.

\subsection{Adaptive Pressure-Aware Compression}

Inspired by research on adaptive compression rate allocation~\cite{adaptive2025,accrag2025}, sqz dynamically adjusts compression intensity based on session pressure---the rate of token injection over a sliding window:

\begin{itemize}
  \item \textbf{Low pressure} ($<$50\% of estimated budget consumed in 30 minutes): Default mode, preserve maximum detail
  \item \textbf{High pressure} ($>$80\%): Escalate to aggressive compression
  \item \textbf{Critical} ($>$90\%): Maximum compression + proactive dedup
\end{itemize}

This implements the core finding from ACC-RAG~\cite{accrag2025} that static compression rates are suboptimal: early in a session detail preservation matters more, while late in a session token scarcity demands harder compression.

\subsection{Entropy-Based Safety Routing}

Not all content should be compressed. Stack traces contain exact line numbers that compression could corrupt. Configuration files with credentials should not be parsed and reformatted. Migration SQL must be preserved byte-exact.

sqz uses Shannon entropy analysis combined with pattern detection to classify content into risk tiers:

\begin{itemize}
  \item \textbf{Safe mode} (0\% compression): Stack traces, private keys, credentials, database migrations
  \item \textbf{Default mode} (20--40\% compression): Normal code, git output, build logs
  \item \textbf{Aggressive mode} (40--60\% compression): Repetitive logs, verbose JSON, boilerplate
\end{itemize}

This addresses the concern raised in~\cite{goalgrad2025} that compression methods which rely on generic importance metrics can inadvertently remove functionally critical tokens.

Version 1.4.0 extended the safety model with three fidelity guarantees derived from field reports:

\begin{itemize}
  \item \textbf{Byte-faithful file reads.} The MCP file tools (\texttt{sqz\_read\_file}, \texttt{sqz\_grep}, \texttt{sqz\_list\_dir}) bypass every content-altering pipeline stage---a file read returns its bytes verbatim apart from ANSI stripping. Token savings on this path come exclusively from dedup references on repeat reads, never from dropping content. (Previously, entropy truncation could silently remove below-median-entropy segments from source files.)
  \item \textbf{Refuse-to-guess formatting.} Structural formatters bail out to raw passthrough rather than emit a partial or misclassified summary. The \texttt{git status} formatter, for example, parses porcelain output strictly, matches localized long-format section headers against git's own translation files (it/de/fr/es/pt\_PT), and passes any unrecognized locale or section through unmodified---a wrong summary is strictly worse than no compression.
  \item \textbf{UTF-8 boundary safety.} All truncation and splitting of user-controlled text rounds to character boundaries, so densely non-ASCII output (Cyrillic, CJK, Arabic, emoji) can never crash the compressor mid-pipeline; the property is enforced by randomized tests over arbitrary strings and cut points.
\end{itemize}

\subsection{N-gram Abbreviation}

For session-level compression beyond dedup, sqz observes recurring multi-token phrases across the session and introduces compact abbreviations. This technique applies BPE-style (Byte Pair Encoding) merging at the output level rather than the tokenizer level---frequently co-occurring phrases are replaced with shorter representations.

Abbreviation is deliberately \emph{lossy for identifiers}: only the first occurrence of a repeated phrase survives verbatim, so a commit SHA or file path embedded in a repeated phrase would be replaced by its abbreviation symbol on later lines. Since 1.4.0 the technique is opt-out (\texttt{--no-abbrev} or \texttt{SQZ\_NO\_ABBREV=1}, mirroring the dedup opt-out) and the generated agent guidance documents the escape hatch, so an agent that needs to copy-paste identifiers verbatim can disable it per command.

% ============================================================
\section{Integration Model}

Unlike prompt compression tools that require API interception or model-specific adapters, sqz integrates at the \textbf{tool execution boundary}---a universal hook point that exists in every agent framework:

\begin{table}[h]
\centering
\small
\begin{tabular}{@{}ll@{}}
\toprule
\textbf{Framework} & \textbf{Hook Mechanism} \\
\midrule
Claude Code & PreToolUse JSON hook \\
Cursor & .cursor/rules/*.mdc \\
Windsurf & .windsurfrules guidance \\
Cline / Roo Code & .clinerules guidance \\
Kiro & .kiro/hooks/ JSON schema \\
Gemini CLI & BeforeTool hook \\
OpenCode & TypeScript plugin \\
Codex & AGENTS.md + MCP entry \\
Zed & .rules + MCP entry \\
Copilot CLI & preToolUse hook \\
Any MCP server & \texttt{sqz-mcp proxy} wrapper \\
\bottomrule
\end{tabular}
\caption{Integration mechanisms across agent frameworks.}
\label{tab:integration}
\end{table}

A single \texttt{sqz init} command detects installed frameworks and configures the appropriate hooks. The compression is transparent---the agent receives compressed output without knowing sqz exists.

Since 1.6.0 the MCP proxy provides a framework-independent integration surface: \texttt{sqz-mcp proxy -- <upstream command>} wraps any stdio MCP server, compresses its tool results through the full pipeline (dedup references on repeats, safe mode for stack traces and secrets, error results untouched), optionally compacts verbose tool descriptions in \texttt{tools/list}, and injects an \texttt{sqz\_expand} tool so the agent can recover any original byte-exact. The client sees a normal MCP server, so this works with every MCP client regardless of whether it supports hooks.

% ============================================================
\section{Evaluation}

\subsection{Token Savings}

Measured across a real developer's week of agentic coding (3,003 compressions):

\begin{table}[h]
\centering
\small
\begin{tabular}{@{}lr@{}}
\toprule
\textbf{Metric} & \textbf{Value} \\
\midrule
Total compressions & 3,003 \\
Tokens in (original) & 721,840 \\
Tokens out (compressed) & 543,398 \\
\textbf{Tokens saved} & \textbf{178,442} \\
Average reduction & 24.7\% \\
Peak reduction (dedup) & 92\% \\
\bottomrule
\end{tabular}
\caption{Aggregate compression statistics from production usage.}
\label{tab:stats}
\end{table}

\subsection{Per-Category Performance}

\begin{table}[h]
\centering
\small
\begin{tabular}{@{}lrrr@{}}
\toprule
\textbf{Content Type} & \textbf{Before} & \textbf{After} & \textbf{Reduction} \\
\midrule
Repeated file reads (5$\times$) & 10,000 & 826 & 92\% \\
Test-fix-test cycle (3 runs) & 15,000 & 5,186 & 65\% \\
Repeated JSON response (3$\times$) & 192 & 79 & 59\% \\
Repeated log lines & 148 & 62 & 58\% \\
Large JSON array & 259 & 142 & 45\% \\
Git diff & 61 & 54 & 12\% \\
Stack trace (safe mode) & 82 & 82 & 0\% \\
\bottomrule
\end{tabular}
\caption{Per-category compression performance.}
\label{tab:percategory}
\end{table}

\subsection{Prompt Cache Impact}

Because sqz operates pre-injection (before content enters the context), the message prefix used for prompt caching remains unchanged. This is a critical advantage over post-hoc pruning methods which modify history and invalidate cached prefixes. Empirically, prompt cache hit rates remain at baseline (90\%+) with sqz, compared to approximately 85\% with history-modifying approaches~\cite{dcp2025}.

\subsection{Latency}

sqz is implemented as a single Rust binary with zero runtime dependencies. Compression latency is sub-millisecond for typical outputs (measured p99 $<$ 5ms for outputs up to 100KB). The SQLite dedup lookup adds approximately 0.1ms per call. Total overhead is negligible compared to LLM inference latency.

% ============================================================
\section{Related Work}

\textbf{Prompt Compression.} LLMLingua~\cite{jiang2023llmlingua} and LLMLingua-2~\cite{pan2024llmlingua2} achieve 2--20$\times$ compression by removing non-essential tokens using a small classifier model. These operate on the \emph{entire prompt} and require an additional LLM inference pass. sqz differs by targeting only tool outputs, requiring no LLM call, and operating at the structural rather than token level.

\textbf{Context Pruning.} ACON~\cite{acon2024} optimizes compression guidelines for long-horizon agents by analyzing failure trajectories. Dynamic Context Pruning~\cite{anagnostidis2023dynamic} manages conversation history post-hoc. These complement sqz---they handle what's \emph{already in context}, while sqz reduces what \emph{enters} context.

\textbf{Structural Compression for Code.} CodeComp~\cite{chen2025codecomp} demonstrates that attention-only compression discards structurally critical tokens (call sites, branch conditions). sqz's domain-specific formatters address this by understanding output structure rather than relying on statistical importance.

\textbf{Information-Theoretic Perspectives.} \cite{infotheory2024} frames compressor-predictor systems through mutual information, showing that larger compressors are more token-efficient. sqz is consistent with this framework---it functions as a deterministic, zero-parameter ``compressor'' that maximizes information per token through structural awareness.

\textbf{Structurally Lossless Trimming.} \cite{dag2025} introduces a three-pass algorithm that strips mechanical bloat (raw tool outputs, base64 images, metadata) while preserving user messages verbatim, achieving up to 86\% reduction. sqz operates at a complementary point---before injection rather than during trimming---but shares the insight that tool outputs are the dominant source of context bloat.

% ============================================================
\section{Limitations and Future Work}

\begin{itemize}
  \item \textbf{Token counting}: the compression pipeline counts tokens with real BPE tokenizers (tiktoken \texttt{o200k\_base} for OpenAI models---exact; \texttt{cl100k\_base} as a $\sim$5\%-variance approximation for Claude and Gemini; chars/4 for unknown local models). The dedup fast path and MCP-side estimates still use the chars/4 heuristic for latency, so aggregate savings figures mix exact and approximate counts.
  \item \textbf{No semantic compression}: sqz does not use an LLM to summarize content. This is by design (zero latency, zero cost, deterministic), but means it cannot capture semantic redundancy across structurally different outputs.
  \item \textbf{Dedup granularity}: exact whole-content hashing, line-aligned slice containment (1.8.0), and near-duplicate deltas via MinHash LSH~\cite{minhash2025} (1.5.0) cover identical, ranged, and lightly-edited re-reads. Overlaps that are not line-aligned, and structural duplicates across different serializations (the same record fetched once as JSON and once as a table), still re-send in full.
\end{itemize}

Future directions include structural dedup across serializations, integration with KV cache compression methods~\cite{chen2025codecomp}, learned formatters that adapt to project-specific output patterns, and an end-to-end evaluation measuring agent task success with and without compression rather than token counts alone.

% ============================================================
\section{Conclusion}

sqz demonstrates that significant token savings (24.7\% average, 92\% on dedup hits) are achievable through pre-injection structural compression---without model fine-tuning, additional LLM calls, or changes to the agent workflow. By operating at the tool output boundary, sqz is framework-agnostic, cache-friendly, and transparent. The combination of domain-specific formatters, content deduplication, and adaptive pressure-awareness produces a practical system that reduces costs and improves session longevity for real-world agentic coding workflows.

\medskip
\noindent\textbf{Availability.} sqz is open source under the Elastic License 2.0. Source code, benchmarks, and documentation are available at \url{https://github.com/ojuschugh1/sqz}.

% ============================================================
\bibliographystyle{plainnat}
\bibliography{references}

\end{document}
